\documentclass[11pt,a4paper]{ctexart}
\usepackage[color]{tlnote}
\usepackage{needspace}
\title{高数 I（A 类）期中 · 考前 24h 精简计划}
\author{陶理}
\noteheader{减量版 · 取代 11/10 执行卡}
\docdate{2026-11-10}

\begin{document}
\renewcommand{\arraystretch}{1.16}
\noindent{\Large\bfseries 考前 24h 精简计划（减量版）}\hfill{\small\color{tlblue}11/10 14:30 起 / 11/11 周三考试}
\par\noindent\rule{\linewidth}{0.4pt}\par\smallskip
\noindent{\small\color{gray}本卡\textbf{取代}《11/10 执行卡》，以本卡为准。旧卡的 108 题安排作废。原则从「加题」改为「抽样式确认 + 保手感 + 早睡」：\textbf{今天新题从 108 压到 24，明天 0 新题}。M1 单独改走《作图首遍训练卡》：先完整做一遍，再谈熟练。}

\section{一、盘点结论：哪里该补、哪里该砍}

\begin{center}\footnotesize
\begin{tabular}{@{}p{2.0cm}p{6.6cm}p{6.4cm}@{}}
\toprule
\textbf{问题} & \textbf{核实到的实数} & \textbf{本卡的处理} \\
\midrule
重复造题 & 六个薄弱点有两套题：定向加练 84（M 段 72）＋ 网络题源 48 ＝ 132 题；库内总题池 301 题 & 今天只取 \textbf{5 题抽检 ＋ M1 首遍完整作图}，其余全部封存 \\
孤儿材料 & 网络题源 48 题已入库、但执行卡一次未提，成了「还欠着」的账 & 明确判定\textbf{考前不做}，考后复盘用 \\
负荷过载 & 11/9 排 184 题；11/10 排 108 题（C18＋B18＋M72），还要留 23:00 停笔 & 新题 24 题（B 卷 18 ＋ 作图 1 ＋ 抽检 5），减 78\% \\
零回执 & 全库 \textbf{0 条做题回执}；所谓「实测薄弱点」实为「自述 × 命题权重」的合成排序；错题回炉无数据可依 & 补一张 \textbf{6 行回执表}（§4），是今天唯一必须写的输出 \\
入口过载 & Downloads 共 39 份 PDF；《讲义\_该翻哪一本》停在 11/04，未覆盖 11/9--11/10 新产出的 6 份 & §5 收口到 \textbf{只碰 3 本} \\
重复任务 & 旧卡要求「出一张考场时间分配表」，但《考前知识点清单》§1 已有 120 min 逐题预算表 & 取消，直接沿用现成表 \\
\bottomrule
\end{tabular}
\end{center}

\noindent{\footnotesize\color{gray} 另注：《高数I期中冲刺计划\_2026-11.md》抬头仍写「9 天冲刺 11/1--11/9、考试约 11/10」（文内已有 11/07 更正行）。属历史文档，本次不重排，只在此标注，避免临考前去翻过期排期。}

\section{二、封存清单（考前明确不做）}

\begin{warnbox}[以下六项一律封存，不做、不补、不算欠账]
\begin{enumerate}[leftmargin=1.8em,itemsep=1pt]
\item \textbf{网络题源 48 题}：整册封存（与定向册同点重复），考后再用。
\item \textbf{定向加练 84 题}：只做 M1 首遍训练 ＋ 五模块各 1 题，其余封存；取消「A 级 56 题」。
\item \textbf{精编本 151 题 / 分享版例题}：只查一句，不刷、不从头翻。
\item \textbf{模拟卷 A、C 全卷}：不重做；旧卡「C 卷 18 题回炉」取消（仅段 1 的分支允许 C 卷第 11 题单题）。
\item \textbf{合订本 370 页}：仅详解查不到时按页码翻一次。
\item \textbf{题池 301}：\textbf{不再追加任何新题 / 新题册}。
\end{enumerate}
\end{warnbox}

\section{三、今天剩余时间（14:30 → 23:00 停笔）}

\noindent{\footnotesize 按当前时点顺排；\textbf{若某段已过，从当前位置接着往下走，绝不回补前面}。休息与睡眠不可被挤占。}

\begin{center}\footnotesize
\begin{tabular}{@{}p{1.0cm}p{1.5cm}p{9.0cm}p{1.2cm}c@{}}
\toprule
\textbf{段} & \textbf{时长} & \textbf{只做这一件事} & \textbf{新题} & \textbf{完成} \\
\midrule
0 & 20 min & 收口：只看《考前知识点清单》§1 时间预算表 ＋ §2 未定式决策表，\textbf{不动笔} & 0 & $\square$ \\
1 & 120 min & \textbf{模拟卷 B 严格掐表}（唯一保留的整卷，保手感）。若今天已完成整卷，改做「C 卷第 11 题 ＋ 3 道今日错题」，\textbf{不重做整卷} & 18 & $\square$ \\
2 & 50 min & 订正整卷：只写错因关键词（展错阶 / 符号反 / 前提漏 / 放缩反），不拓展同型 & 0 & $\square$ \\
3 & 50 min & 离开桌子，吃饭走动 & 0 & $\square$ \\
4 & 60 min & \textbf{M1 首遍完整作图 50 min}（《作图首遍训练卡》：照抄 8 ＋ 半独立 15 ＋ 独立 22 ＋ 对答案 5）＋ \textbf{其余五模块各 1 题抽检 10 min} & 6 & $\square$ \\
5 & 40 min & 错题回炉：只重做今天标错的题，每道只重写到卡住那一步，不整题重抄 & 0 & $\square$ \\
6 & 40 min & 《考前知识点清单》末页 \textbf{14 项闭卷默写}，卡住的回对应小节看 3 min & 0 & $\square$ \\
7 & 30 min & 《泰勒展开全表》只看 §2 核心九个 ＋ §8 六个陷阱，其余页不翻 & 0 & $\square$ \\
8 & 20 min & 填 §4 的 \textbf{6 行回执}——今天唯一必须留下的输出 & 0 & $\square$ \\
9 & 40 min & 考场物品 / 确认 11/11 时间线 / 放松；\textbf{无加码条款} & 0 & $\square$ \\
\midrule
\multicolumn{3}{r}{合计新题} & \textbf{24} & 旧卡 108 \\
\bottomrule
\end{tabular}
\end{center}

\begin{keypoint}[验收线与放弃线]
\begin{itemize}\itemsep=1pt
\item \textbf{验收线}：整卷正确率 $\ge 80\%$；五模块抽检各 $1/1$；M1 按训练卡 §6 七项自检 $\ge 6$ 项过；默写 14 项 $\ge 12$ 项过。
\item \textbf{降级规则}：某模块抽检错，\textbf{只回炉 2 道同型}；M1 自检不过只重画那张表，\textbf{不补第二道作图题}。
\item \textbf{放弃线}：单题卡 $>8$ min 直接跳；任一段超时就砍掉后面最不重要的一段，\textbf{不挤占段 3 休息与睡眠}。
\item \textbf{23:00 停笔 / 23:30 睡}——不压缩，旧卡这条保留。
\end{itemize}
\end{keypoint}

\section{四、M1 首遍完整作图 ＋ 五模块抽检 5 题 ＋ 6 行回执}

\noindent{\footnotesize \textbf{M1 单独处理}：你自述「还没完整做过一遍画图、列表不熟」。核查确认两套 M1 题共 22 道\textbf{全是零件}（渐近线 / 驻点 / 拐点），无一道要求列表或画图。故 M1 不走抽检，改走《作图首遍训练卡\_11-10.pdf》三遍法（50 min）：照抄五行表 $\to$ 半独立填表 $\to$ 独立做 1 道完整作图（22 min 掐表）。完成 $\square$　七项自检通过 \blank[1.5cm] / 7。}

\noindent{\footnotesize 其余五模块各取\textbf{基础第 1 题} 1 道（题号见 \texttt{讲义\_薄弱点定向加练\_11-09.pdf}），只用于「确认没退化」：}

\begin{center}\footnotesize
\begin{tabular}{@{}p{1.2cm}p{4.2cm}p{3.0cm}p{2.0cm}p{1.4cm}p{2.0cm}@{}}
\toprule
\textbf{模块} & \textbf{主题} & \textbf{取题} & \textbf{对 / 1} & \textbf{过?} & \textbf{错因关键词} \\
\midrule
M2 & 量词行 / M 型 / 放缩 & M2.1 & \blank[1.5cm] & $\square$ & \blank[1.6cm] \\
M3 & 条件书写 / 六型求导 & M3.1 & \blank[1.5cm] & $\square$ & \blank[1.6cm] \\
M4 & 泰勒：阶 / 系数 / 符号 & M4.1 & \blank[1.5cm] & $\square$ & \blank[1.6cm] \\
M5 & 极限计算 / 间断点 & M5.1 & \blank[1.5cm] & $\square$ & \blank[1.6cm] \\
M6 & 新压轴 / 极坐标 & M6.1 & \blank[1.5cm] & $\square$ & \blank[1.6cm] \\
\bottomrule
\end{tabular}
\end{center}

\needspace{9\baselineskip}
\vspace{2pt}
\noindent{\footnotesize\textbf{6 行回执}（只填数字与一句话，不写复盘。这是把「合成排序」升级为「实测排序」的唯一办法）：}

\begin{center}\footnotesize
\begin{tabular}{@{}p{4.6cm}p{2.2cm}p{7.4cm}@{}}
\toprule
\textbf{项目（分母已标）} & \textbf{结果} & \textbf{最该带走的一句话} \\
\midrule
整卷正确率 / 18 & \blank[1.6cm] & \blank[7.0cm] \\
抽检合计 / 5 & \blank[1.6cm] & \blank[7.0cm] \\
默写自测 / 14 & \blank[1.6cm] & \blank[7.0cm] \\
最卡的一个模块 & \blank[1.6cm] & \blank[7.0cm] \\
明天只回看哪 2 处 & \blank[1.6cm] & \blank[7.0cm] \\
停笔时间 & \blank[1.6cm] & \blank[7.0cm] \\
\bottomrule
\end{tabular}
\end{center}

\section{五、11/11 考前 40 min（0 新题）}

\begin{keypoint}[资料收口：只碰 3 本 ＋ 1 张表]
\textbf{《考前知识点清单》7 页}（主册：§1 时间预算、§2 决策表、§7.3 作图七步、§9 错误清单、末页默写）$\mid$
\textbf{《泰勒展开全表》4 页}（只 §2 ＋ §8）$\mid$
\textbf{《模拟卷 B》}（只在需要确认某题型手感时翻一题，不重做）。
\textbf{＋《作图首遍训练卡》§5 标准答案表}（明早只 30 秒过一眼五行结构，不重做）。其余 PDF 考前一律不打开。
\end{keypoint}

\begin{center}\footnotesize
\begin{tabular}{@{}p{1.6cm}p{12.6cm}@{}}
\toprule
\textbf{时长} & \textbf{动作} \\
\midrule
15 min & 知识点清单 §1 逐题时间预算 ＋ §9 十五条高频错误 ＋ 止损三条 \\
10 min & 把 §2 未定式决策表与 §7.3 作图七步表\textbf{默到草稿纸角上}（进考场后第一件事） \\
10 min & 泰勒全表 §2 核心九个展开，只过一遍符号与阶 \\
5 min & 末页 14 项默写自测：\textbf{只勾不熟的}，熟的直接跳过 \\
\bottomrule
\end{tabular}
\end{center}

\noindent{\footnotesize\textbf{考场止损三条}：卡 5 分钟就走（选择 3 分钟先猜后标）$\mid$ 不会也写开头（量词行、作图七步表先落纸）$\mid$ 留 5 分钟只查符号、端点、定义域。}

\vspace{8pt}
\noindent{\footnotesize\color{gray} 本卡为独立 PDF，不改动任何既有讲义 / 合订本 / 阅读包的绝对页码；旧《11/10 执行卡》保留但不执行。}

\end{document}
